\section{Discussion: interpretation and implications}
\label{sec:discussion}

\subsection{What the estimate measures}

The estimate is the share of measured value added produced by informal work, as defined by the labor classification of the household survey. Proposition~\ref{prop:share} links it to three measurable quantities: the informality rate, relative earnings and the formal capital share. The link uses the survey concepts directly. The income of own-account workers and informal employers is mixed income, which already contains the return to their own capital, and the income of informal employees is adjusted by group (Section~\ref{sec:data}). The posterior bands report uncertainty given this mapping, and Table~\ref{tab:sensitivity} reports how the level moves under alternative mappings.

This concept is related to, but not the same as, the shadow or non-observed economy. Formal firms may employ workers informally, some own-account production is declared and included in the accounts, and illegal production is outside the labor classification used here. The comparisons in Figure~\ref{fig:benchmarks} are therefore checks of magnitude and trend. They are reassuring on both counts: the estimate lies within the range of published series, very close to the enhanced MIMIC estimates in 2009--2020, and shares their downward trend.

\subsection{Identifying the path}

The identification exercises separate two objects. The steady-state parameter $\theta$ is not pinned down by aggregate data, which explains why calibrated and estimated DGE models in the literature can reach very different levels for the same country. The quarterly path of the informal share, which is the object of interest, is identified once relative earnings are observed: from 2007 it is pinned down by the observed informality rate, the observed earnings ratio and the capital share, and it is invariant to $\theta$ within 0.3 percentage points. Before 2007 the path is identified with the information available for those years, the reconstructed informality rate and the national accounts, and the model carries the earnings information back; this step moves the 1990s by at most 1.4 points across values of $\theta$ between 0.10 and 0.40. Simulations with known parameters recover the true share with high precision.

The practical implication for this literature is that the level of the informal economy in a DGE model should be identified by a direct measure of relative productivity rather than by the aggregates or by calibration, and that the identification check should be reported together with the posterior.

\subsection{Implications for research and policy}

The estimates show that the contribution of informal work to output has declined together with labor informality, from about 30\% to 20\% of GDP over three decades. They also quantify the productivity gap behind the difference between the two measures: at the end of the sample informal value added per worker is about one-fifth of formal value added per worker. This gap is the relevant magnitude for policy discussions. Formalizing a worker does not by itself raise output; what matters is whether the worker moves into higher-productivity activity. Evaluating specific reforms, such as changes in payroll taxes or enforcement, is a natural next step for the model, which already contains the formality wedge through which such policies operate.

Because the data and code are public and the earnings series is built from the regular household survey, the estimate can be updated every quarter. It thus offers a timely complement to DANE's periodic measurement of the non-observed economy.
